% Source: https://artofproblemsolving.com/community/c3h2615498_peace09s_3000th_post
% Stable identity: aops-topic-2615498-q16
Pair the factor $i$ with $n-i$: $(x+i)(x+n-i)=x^2+nx+i(n-i)=i(n-i)-y$. Since $n$ is odd, the $n+1$ factors split into the $\dfrac{n+1}2$ pairs $i=0,1,\dots,\dfrac{n-1}2$. So
\[\prod_{i=0}^n(x+i)^2=\prod_{i=0}^{\dfrac{n-1}2}\bigl(y-c_i\bigr)^2,\qquad c_i=i(n-i).\]
This is a monic polynomial in $y$ of degree $n+1$.

\textbf{Coefficient $a_n$.} It is $-2\sum c_i$. Summing $i(n-i)$ over all $i=0,\dots,n$ gives $n\cdot\dfrac{n(n+1)}2-\dfrac{n(n+1)(2n+1)}6=\dfrac{n(n+1)(n-1)}6$, and by the symmetry $i\leftrightarrow n-i$ (no middle term, as $n$ is odd) the sum over $i\le\dfrac{n-1}2$ is half of this. So $a_n=-2\cdot\dfrac{n^3-n}{12}=\dfrac{n-n^3}6$.

\textbf{Coefficient $a_2$.} Since $c_0=0$, the product is $y^2\prod_{i\ge1}(y-c_i)^2$, so $a_2=\prod_{i=1}^{\dfrac{n-1}2}c_i^2$. And $\prod_{i=1}^{\dfrac{n-1}2}i(n-i)=1\cdot2\cdots\dfrac{n-1}2\cdot\dfrac{n+1}2\cdots(n-1)=(n-1)!$. So $a_2=\bigl((n-1)!\bigr)^2$. $\blacksquare$
